\subsection{Single-asset generator validation}

We first evaluated whether the asset generator captured the return
distribution and volatility clustering that motivated its use in
composition (Figs.~\ref{fig:empirical_motivation} and
\ref{fig:model_architecture}; Table~\ref{tab:descriptive_stats}). The observed SPY series had heavy tails
and persistent absolute-return autocorrelation in both the training
and holdout windows. Individual raw-return correlations were small,
although the Ljung--Box test detected joint serial dependence
(Table~\ref{tab:descriptive_stats}). Calibration against the
$2014$--$2024$ history selected jump probability and mean duration
$(\epsilon,\lambda)=(10^{-4},90)$, with similar objective values at
nearby grid points (Eq.~\eqref{eq:grid_search};
Table~\ref{tab:hyperparams}; Algorithm~\ref{alg:hmmwj}). Ordinary
states lasted only $1$--$2$ steps, and the counted transition matrix
implied absolute-return autocorrelation that decayed to near zero
within a week (Fig.~\ref{fig:model_internals}). The selected episodes
extended visits to extreme states, providing a mechanism for more
persistent volatility.
To assess whether the generator comparison depended on the state
partition, we varied the number of states
(Table~\ref{tab:sensitivity_n}). The primary SPY analysis used $N=100$ states, while
sweeps from $N=30$ to $200$ changed HMM-NJ pass rates by less than
one percentage point (Table~\ref{tab:sensitivity_n}). The HMM-WJ
model showed a tradeoff as the partition changed: narrower states
improved autocorrelation accuracy but lowered the AD pass rate.
States with too few observations use the documented global-moment
fallback instead of local emission estimates
(Table~\ref{tab:fallback_assets}). The selected
$N=100$ lay within the range that supported well-populated states
for the $2{,}766$-day history. That choice did not establish an
appropriate resolution for shorter histories, which would provide
fewer observations per state. The sensitivity checks therefore
supported the fitted generator at the observed sample length without
establishing one state count for every asset history.

% Main table tab:model_comparison is collected after the references.

We compared the calibrated HMM-WJ with seven alternative generators
using $1{,}000$ paths of $2{,}766$ trading days
(Table~\ref{tab:model_comparison}; Fig.~\ref{fig:model_comparison}).
The HMM-NJ and HMM-WJ models had KS pass rates of $99.3\%$ and
$98.3\%$, respectively, compared with $98.4\%$ for the fitted
Laplace baseline. Both HMM variants had lower Wasserstein-1 distances
than Laplace, whereas Laplace had the lowest Hellinger distance among
the fitted generators (Table~\ref{tab:model_comparison}). The
ranking therefore depended on the metric. Bootstrap resampling had a
$100\%$ KS pass rate but did not reproduce volatility clustering,
while Gaussian sampling passed neither distribution test.
The GARCH(1,1) model had the lowest mean absolute error in the
autocorrelation function (ACF-MAE), $0.031$, but rarely passed the
distribution tests. The gated recurrent unit (GRU)~\cite{choEtAl2014gru}
showed a similar tradeoff. The HSMM-style baseline adapted from Bulla
and Bulla~\cite{bullaBulla2006} reached an $82\%$ KS pass rate but
underestimated kurtosis by more than one third. Within the HMM,
Student-$t$ emissions produced kurtosis of $7.5$, close to the
observed $7.7$, whereas Normal emissions underestimated it by about
$29\%$ (Table~\ref{tab:emission_hsmm}). Both HMM variants also
reproduced the observed skewness of $-0.75$ through unequal occupancy
of negative and positive states (Table~\ref{tab:descriptive_stats}).
The temporal comparison showed the cost of retaining that marginal fit.
Jump episodes improved temporal fit, with a small reduction in
distributional pass rates (Table~\ref{tab:model_comparison};
Fig.~\ref{fig:model_comparison}b). The HMM-WJ model reduced ACF-MAE from
$0.059$ for HMM-NJ to $0.053$. About $24\%$ of paths contained an
episode (Fig.~\ref{fig:statistical_validation}); paths without an
episode were equivalent to HMM-NJ. The
episode share, the fraction of forced steps, and the jump-path
autocorrelation agreed with the calculations in Online
Appendix~\ref{sec:supp-generator}
(Tables~\ref{tab:supp-generator-quantities} and
\ref{tab:supp-acf-prediction}). To examine the tradeoff, we
resampled the fitted ensemble with the fraction of jump-active paths
ranging from zero to one. Increasing that fraction lowered
autocorrelation error but also lowered KS and AD pass rates.
Kurtosis declined through the observed value, while the Hill
upper-tail index remained stable. The fitted ensemble retained approximately one quarter jump-active
paths, balancing kurtosis and autocorrelation error.
To test whether the distributional fit extended beyond calibration,
we compared frozen-model paths with the $249$ held-out SPY returns
from $2025$. The HMM-NJ and HMM-WJ models retained KS pass rates of
$96.6\%$ and $95.4\%$, with kurtosis close to the observed value
(Table~\ref{tab:model_comparison};
Fig.~\ref{fig:statistical_validation}). The fitted Laplace
baseline had a KS pass rate of $96.7\%$ and a lower Wasserstein-1
distance than either HMM variant, but its excess kurtosis of $2.8$
remained below the observed $6.9$ (Table~\ref{tab:model_comparison}).
Separate calibrations for
NVDA, JNJ, and JPM also showed strong training fit but more variable
holdout performance (Table~\ref{tab:cross_asset}). These findings
supported the HMM as a source of heavy-tailed asset paths, while the
benefit of the jump mechanism depended on the metric and evaluation
period. We used the no-jump variant for the primary composition
comparison and evaluated the transfer of jump settings separately
in the multi-asset jump comparison.

% Main figure fig:model_comparison is collected after the references.

\subsection{Variance-corrected multi-asset composition}

We compared six composition methods across a training universe of
$423$ non-market assets, with $393$ eligible for GARCH, to test
whether the correction preserved generator variance and
calibrated market loading (Table~\ref{tab:aggregate};
Fig.~\ref{fig:preservation}). We paired each composed path with its
uncomposed generator draw to measure the variance ratio directly.
Across the $421$ ordinary assets, the median of the per-ticker
median ratios was $1.3171$ for naive composition and $1.0001$ for
hybrid composition (Fig.~\ref{fig:preservation}b). The corresponding
hybrid ticker medians ranged from $0.9933$ to $1.0072$. Hybrid also
retained higher marginal-test pass fractions across the calibrated
$\beta$ range (Fig.~\ref{fig:preservation}a). The two tracker assets
used a different variance target and were identified separately.
Finite-path covariance still caused variation around the ordinary
hybrid target: the median absolute relative variance error was
$1.12\%$, and its $95$th percentile was $3.41\%$
(Fig.~\ref{fig:preservation}b). These path-level errors qualified the
small average cross-covariance reported in the earlier diagnostic
(Fig.~\ref{fig:cov_diag}): cancellation across replications did not imply
exact variance preservation on each path.
The correction improved distributional fit while recovering the
market loading with accuracy similar to the other methods
(Table~\ref{tab:aggregate}).
All six methods had median absolute $\beta$ errors between $0.017$
and $0.022$ (Table~\ref{tab:aggregate}). Centering also kept the
median absolute intercept error at $0.002\,\mathrm{yr}^{-1}$ for
the naive and corrected methods; residual-fit methods retained
finite-path variation in their residual means. With observed training
SPY and no loading adjustment, centered paths matched the observed
asset mean exactly. Their small intercept errors therefore reflected
the imposed mean constraint and finite-path loading error, rather
than independent evidence of mean recovery.

% Main figure fig:preservation is collected after the references.

% Main table tab:aggregate is collected after the references.

To test whether fallback emissions influenced the composition
comparison, we inspected the fitted states and ran a paired
sensitivity check (Tables~\ref{tab:fallback_assets} and~\ref{tab:fallback_sensitivity}). Fallbacks affected $50$ states in $39$ full-return
models, including $39$ populated states containing repeated returns;
no residual-fit model required a fallback. Keeping the empirical
spread in those populated states raised hybrid KS pass rates from
$68.6\%$ to $78.9\%$ and AD pass rates from $56.2\%$ to $70.2\%$
on the same $39$ affected assets (Table~\ref{tab:fallback_sensitivity}). The alternative retained
partitions, transition matrices, state paths, and state means.
Treating repeated returns as noise on the scale of the entire asset
therefore affected the reported training fit. The main comparisons
retain the original fallback rule; the control identifies sensitivity
to that rule without establishing the alternative's holdout behavior.

The methods differed more in distributional fit than in factor
recovery (Table~\ref{tab:aggregate}; Fig.~\ref{fig:tails}).
Gaussian SIM had the lowest median $R^2$ error, but almost never
passed the distribution tests and had excess kurtosis an order of
magnitude below the observed value. The corrected method passed $73.6\%$ of KS
comparisons and retained heavy tails, with median excess kurtosis
of $7.2$ and a Hill index of $0.37$ (Fig.~\ref{fig:tails}). Its
median $R^2$ error fell between those of the naive and Gaussian
methods because the main correction targeted generator variance
rather than observed $R^2$. Preserving variance therefore improved
the marginal fit without making every fitted property an exact
target.
Fitting directly to regression residuals provided a competitive
alternative because those residuals already excluded the market
component (Table~\ref{tab:aggregate}). The three residual-fit methods had KS pass rates ranging
from $67.4\%$ for GARCH(1,1)-$t$ to $75.8\%$ for
JumpHMM-on-residuals (Table~\ref{tab:aggregate}). The corrected rate
was about two percentage points below JumpHMM-on-residuals and
slightly above block bootstrap. JumpHMM-on-residuals and block
bootstrap had slightly lower Wasserstein distances, while
JumpHMM-on-residuals matched the corrected method's kurtosis and Hill
index almost exactly. The block-bootstrap and GARCH residual fits
produced lighter tails. These comparisons used each method's native
mean treatment, so their cross-method differences also included the
effect of the zero-sum constraint.
Applying equivalent residual centering to all six methods changed the
training ranking on the common $393$-asset set
(Table~\ref{tab:centering_control}). Hybrid retained a $74.4\%$ KS
pass rate, while JumpHMM-on-residuals rose from $75.8\%$ to $78.4\%$
and block bootstrap rose from $73.5\%$ to $76.7\%$. Both residual-fit
methods also had higher AD pass rates and lower Wasserstein distances
than hybrid after centering. Gaussian SIM and GARCH remained below
hybrid on both pass rates. After centering, all methods had median
intercept errors of $0.002\,\mathrm{yr}^{-1}$ at the reported precision,
while leaving each path's variance, recovered loading, $R^2$, and
excess kurtosis unchanged to numerical tolerance. Hybrid's improvement
over naive composition persisted because both already used identical
centering. Reuse therefore improved on naive composition without a
second generator fit, but did not match the strongest residual-fit
comparators under the common mean constraint. Held-out histories were
needed to assess whether the improvement over naive composition
persisted beyond the fitting period.

\subsection{Holdout performance and VaR coverage}

To test out-of-sample generalization without refitting or look-ahead
bias, we froze the fitted models and SIM parameters before evaluating
them against $416$ complete asset histories from $2025$
(Table~\ref{tab:oos_composition}). The corrected
method reached an $87.0\%$ mean cross-ticker KS pass rate, against
$83.5\%$ for centered naive composition, and its median composed-to-observed
variance ratio was $0.97$, whereas naive composition remained
inflated at $1.29$. JumpHMM-on-residuals reached an $86.2\%$ KS pass
rate and a slightly lower median $W_1$ than the corrected method. Matched-length
training blocks produced lower pass rates for every method. The higher
$2025$ pass rates did not arise solely from comparing samples of
$249$ rather than $2{,}766$ observations.
The average pass rates concealed substantial differences among
tickers (Fig.~\ref{fig:oos_composition}). Most ticker-level rates
were near the upper end of the range, but every method retained a
lower tail of weak fits. For the corrected method, many tickers improved
relative to their matched-length training blocks, while a subset
deteriorated out of sample.
The held-out marginal improvement was broad but did not ensure
reliable extreme-loss thresholds.
To test whether marginal fit translated into reliable tail thresholds,
we estimated per-ticker one-day left-tail Value-at-Risk (VaR)
thresholds from $5{,}000$ separate $249$-day paths and counted
exceedances on the observed $2025$ histories (Table~\ref{tab:var}).
At $\alpha=0.95$, the corrected and JumpHMM-on-residuals methods
produced mean exceedance rates of $5.48\%$ and $5.62\%$, respectively,
against the $5\%$ nominal rate. At $\alpha=0.99$, their rates were
$1.12\%$ and $1.16\%$, against the $1\%$ nominal rate. Naive
composition was below nominal at $0.88\%$, while Gaussian SIM had
the largest rate at $1.72\%$. Doubling the threshold pool from
$2{,}500$ to $5{,}000$ paths changed every method's mean exceedance
rate by less than $0.006$ percentage points at either coverage level
(Table~\ref{tab:var_convergence}). The corrected method and
naive composition lay similarly close to the $1\%$ target on
opposite sides. Their small difference did not establish a ranking:
the holdout contained only about $2$--$3$ expected $99\%$ exceedances
per asset, and all assets shared one market period. The corrected
method's mean rates remained above nominal at both coverage levels. Cross-asset dispersion described the
heterogeneity of realized coverage; it did not establish statistical
miscalibration or validate coverage in other periods.

% Main table tab:oos_composition is collected after the references.

% Main figure fig:oos_composition is collected after the references.

% Main table tab:var is collected after the references.

\subsection{Jump-enabled multi-asset composition}
\label{sec:results-jump-composition}

To test whether jump episodes improved the composed paths, we compared
jumps disabled, market-only jumps, and market-plus-asset jumps under
naive and corrected composition (Table~\ref{tab:jump_ablation}).
With corrected composition, market-only jumps reduced the training
absolute-return ACF-MAE over lags $1$--$25$ from $0.13569$ to
$0.12844$, a $5.3\%$ reduction. The KS pass rate changed from
$73.94\%$ to $73.64\%$, with a Monte Carlo interval for the difference
that included zero (Table~\ref{tab:jump_ablation_uncertainty}).
Enabling jumps in both the market and asset generators reduced the
ACF-MAE further, to $0.11581$, a $14.6\%$ reduction from the no-jump
setting, but lowered the KS pass rate to $61.68\%$. Across all settings
and both windows, the mean ratio of corrected to generator variance
differed from one by less than $0.05\%$. The correction remained effective,
while asset episodes reduced training autocorrelation error at the
cost of marginal fit.
The training-period temporal improvement did not carry into the
$2025$ holdout (Table~\ref{tab:jump_ablation}). With corrected composition, market-only jumps raised
the $25$-lag ACF-MAE from $0.07330$ to $0.07525$, a $2.7\%$
increase, while market-plus-asset jumps raised it to $0.07832$, a
$6.9\%$ increase (Table~\ref{tab:jump_ablation}). The corresponding
KS pass rates were $85.92\%$, $85.88\%$, and $84.43\%$.
The Monte Carlo intervals excluded zero for both increases in ACF
error, and the $60$-lag results agreed
(Table~\ref{tab:jump_ablation_uncertainty}). The jump-active market
paths represented only $2.2\%$ of the simulated one-year ensemble
(Table~\ref{tab:jump_ablation_episodes}), so the holdout comparison tested the fitted mixture of ordinary and
jump-active paths. The correction worked with jumps enabled, but the transferred SPY
settings did not improve temporal fit in $2025$. The short holdout
and low episode frequency limited what the comparison could establish
about behavior in other market periods.

% Main table tab:jump_ablation is collected after the references.

\subsection{Sensitivity analyses}

We examined the tracker and clipping rules under conditions where
they could affect the result (Figs.~\ref{fig:branch_map} and
\ref{fig:r2_distribution}; Table~\ref{tab:branch}). In the observed universe, $421$ assets
used the variance-preserving branch and two, QQQ and SPYG, used the
$R^2$-preserving branch; none required clipping
(Figs.~\ref{fig:r2_distribution} and \ref{fig:branch_map};
Table~\ref{tab:branch}). Both trackers had $100\%$ KS pass rates,
while the highest-$R^2$ individual stock, BLK, fell about $0.16$
below the tracker threshold of $0.80$. Because only two assets used
that branch, we tested a synthetic-tracker grid with
$(\beta_{\mathrm{true}},R^2_{\mathrm{true}})\in\{0.8,1.0,1.2\}
\times\{0.80,0.85,0.90,0.95,0.99\}$. Each cell had a $100\%$ KS pass
rate and recovered its target loading and $R^2$
(Table~\ref{tab:tracker_grid}). These checks supported the tracker
rule across the tested grid, while its empirical evaluation remained
limited to two assets.
Greater market exposure made the composed distribution more dependent
on the market path (Fig.~\ref{fig:preservation}a;
Table~\ref{tab:beta_bucket}). Grouping assets by $\beta$ quartile showed that
corrected KS pass rates fell from $87.6\%$ in the lowest quartile
to $65.9\%$ in the highest (Table~\ref{tab:beta_bucket}).
To test the loading adjustment directly, we multiplied the market
growth-rate path by $\gamma\in\{1,2,3\}$ while holding the calibrated
marginals fixed (Table~\ref{tab:stress}). At $\gamma=1$, every
asset had $\rho_i<1-f=0.90$, so clipping was unnecessary.
At $\gamma=2$ and $3$, clipping applied to $76.4\%$ and $95.2\%$
of assets. The adjustment therefore activated as intended when the
market contribution would otherwise consume too much of the target
variance. At observed market variance, varying the residual floor
$f$ did not change branch assignments, and varying the tracker
threshold changed only the two tracker assignments
(Table~\ref{tab:branch}).
